\section{Significance of the Study}


This study contributes to the few-shot image generation domain by enhancing the editability and interpretability 
of latent codes in the SAGE framework. By replacing the PSP encoder with the StyleFeatureEditor's encoder, the 
proposed FeatureSAGE framework improves stability and consistency in attribute group manipulation under 
data-scarce cases which addresses the key limitation of the existing few-shot generative models. 

Aside from methodological contributions, the proposed approach also has practical relevance and has several 
downstream applications. In fine-grained visual recognition, FeatureSAGE can generate identity-preserving 
variations of an animal or plant from a very few dataset, which helps train classifiers for species or breed 
identification in environments where collecting large dataset is impossible or difficult to acquire. In biometric 
work that focuses on identity, such as forensic image analysis and identity-based data augmentation, FeatureSAGE 
can generate multiple edited versions of an image while keeping the same person’s identity. This makes the findings 
more realistic and more reliable. In medical and scientific imaging, researchers often have small datasets because 
privacy rules limit sharing, and gathering new images can be expensive. In some cases, the condition or disease is 
rare.FeatureSAGE generates realistic images that can serve as a synthetic dataset for training and validating 
machine learning models. 

This study extends prior work on attribute group editing and GAN latent spaces, and it gives clear guidance for 
building data-efficient image generation systems that can support downstream applications that need consistent, 
reliable image synthesis.